\uuid{0QaI}
\titre{Graphe}
\auteur{bodin ; hanani ; mzari}
\organisation{exo7}
\licence{CC BY-NC-SA 4.0}
\datecreate{2020-01-20}
\niveau{L1}
\module{Analyse}
\chapitre{Fonctions circulaires et hyperboliques inverses}
\sousChapitre{Autre}
\difficulte{2}
\theme{Fonctions usuelles, graphe}
\exo7qcmid{lille-1.fonctions-usuelles.019}

\contenu{
\begin{question}
\qclassification{126.00}
\qtags{motcle=graphe}
Soit $f(x)=\sqrt{\frac{x^3}{x-1}}$.  Quelles sont les assertions vraies ?

\begin{answers}
    \good{Le domaine de définition de $f$ est $]-\infty, 0]\cup ]1,+\infty[$.}

    \bad{$y=x-\frac{1}{2}$ est une asymptote à la courbe de $f$ en $+\infty$.}

    \good{$y=x+\frac{1}{2}$ est une asymptote à la courbe de $f$ en $+\infty$.}

    \good{$y=-x-\frac{1}{2}$ est une asymptote en $-\infty$.}

\end{answers}
\begin{explanations}
$f$ est définie si $x\neq 1$ et $\frac{x}{x-1}\ge 0$. On déduit que le domaine de définition de $f$ est $]-\infty, 0]\cup ]1,+\infty[$.

$\lim_{x \to \pm \infty}f(x) = + \infty$,  $\lim_{x \to + \infty}\frac{f(x)}{x} = 1$, $\lim_{x \to + \infty}(f(x)-x) = \lim_{x \to + \infty}x(\sqrt{\frac{x}{x-1}}-1)= \lim_{x \to + \infty} \frac{x}{(x-1)(\sqrt{\frac{x}{x-1}}+1)}= \frac{1}{2}$. Donc $y=x+\frac{1}{2}$ est une asymptote à la courbe de $f$ en $+\infty$.

$\lim_{x \to - \infty}\frac{f(x)}{x} = -1$,  $\lim_{x \to - \infty}(f(x)+x) = \lim_{x \to - \infty}x(1-\sqrt{\frac{x}{x-1}})= -\frac{1}{2}$. Donc $y=-x-  \frac{1}{2}$ est une asymptote à la courbe de $f$ en $-\infty$.
\end{explanations}
\end{question}
}
